\clearpage\pagestyle{otMain}\fontsize{11}{15}\selectfont\corrigesChapitre{5}{Équations, inéquations et produits nuls}
\begin{blocLivre}{cor-5-1}\begin{corrige}{1}{Une opération inverse}\demarche{Appliquer l’opération inverse aux deux membres.}
On soustrait 13 à la première égalité : $\res{x=-9}$. On divise la seconde par $-4$ : $\res{x=-9/2}$. Les deux équations sont indépendantes.
\end{corrige}\end{blocLivre}
\begin{blocLivre}{cor-5-2}\begin{corrige}{2}{Deux étapes}\demarche{Isoler le produit avant l’inconnue.}
On ajoute 5 : $7x=35$. Puis on divise par 7 : $\res{x=5}$. Contrôle : $7\times5-5=30$.
\end{corrige}\end{blocLivre}
\begin{blocLivre}{cor-5-3}\begin{corrige}{3}{L’inconnue des deux côtés}\demarche{Rassembler les termes inconnus avant les constantes.}
On soustrait $5x$ : $3x+3=-12$. On retire 3 : $3x=-15$. Ainsi $\res{x=-5}$. Les deux membres initiaux valent $-37$.
\end{corrige}\end{blocLivre}
\begin{blocLivre}{cor-5-4}\begin{corrige}{4}{Avec une parenthèse}\demarche{Développer avant de regrouper.}
On développe : $6x-15=4x+9$, donc $2x=24$ et $\res{x=12}$. Vérification : $3(24-5)=57=48+9$.
\end{corrige}\end{blocLivre}
\begin{blocLivre}{cor-5-5}\begin{corrige}{5}{Une première inéquation}\demarche{Regarder le signe du nombre par lequel on divise.}
En divisant par 3, positif, on obtient $\res{x\ge4}$. En divisant par $-2$, négatif, le sens change : $\res{x\le-4}$. Les crochets fermés indiquent que 4 et $-4$ sont solutions.
\label{fig-cor-5-5}\figureGeo[corrige]{inequation-droites}
\end{corrige}\end{blocLivre}
\begin{blocLivre}{cor-5-6}\begin{corrige}{6}{Une équation fractionnaire}\demarche{Multiplier chaque terme par le même dénominateur commun.}
On multiplie toute l’égalité par 12 : $4x+3x=168$. Donc $7x=168$ et $\res{x=24}$. Contrôle : $24/3+24/4=8+6=14$.
\end{corrige}\end{blocLivre}
\begin{blocLivre}{cor-5-7}\begin{corrige}{7}{Tout ou rien}\demarche{Interpréter l’égalité restante lorsque l’inconnue disparaît.}
La première égalité se réduit à $4x-8=4x-8$ : $\res{\text{tous les nombres sont solutions}}$. La seconde imposerait $-8=1$ après soustraction de $4x$ : elle n’a $\res{\text{aucune solution}}$.
\end{corrige}\end{blocLivre}
\begin{blocLivre}{cor-5-8}\begin{corrige}{8}{Un produit nul}\demarche{Séparer les deux facteurs, sans développer.}
Au moins un facteur est nul : $2x-7=0$ ou $x+4=0$. On obtient $\res{x=7/2\text{ ou }x=-4}$. Les deux valeurs annulent bien le produit.
\end{corrige}\end{blocLivre}
\begin{blocLivre}{cor-5-9}\begin{corrige}{9}{Factoriser pour résoudre}\demarche{Factoriser, par un facteur commun ou une identité, avant d’appliquer le produit nul.}
$x^2-6x=x(x-6)$, donc $\res{x=0\text{ ou }x=6}$ ; diviser par $x$ aurait perdu 0. Ensuite $(x+3)^2-4^2=(x+3-4)(x+3+4)=(x-1)(x+7)$ : $\res{x=1\text{ ou }x=-7}$.
\end{corrige}\end{blocLivre}
\begin{blocLivre}{cor-5-10}\begin{corrige}{10}{Lire une inéquation sur un graphique}\demarche{Lire les abscisses des points situés sur ou au-dessus de la droite horizontale.}
La droite horizontale $y=6$ coupe la droite au point d’abscisse 4 ; la droite est au-dessus pour les abscisses plus grandes. Les solutions sont $\res{x\ge4}$. Par le calcul, on divise par 1,5, positif : $x\ge6/1{,}5=4$.
\label{fig-cor-5-10}\figureGeo[corrige]{inequation-graphique}
\end{corrige}\end{blocLivre}
\begin{blocLivre}{cor-5-11}\begin{corrige}{11}{Des billets de spectacle}\demarche{Exprimer les deux effectifs avec une seule inconnue.}
Notons $x$ le nombre de jeunes ; il y a $28-x$ autres billets. L’équation est $6x+10(28-x)=232$, soit $280-4x=232$. Ainsi $\res{x=12}$. Il reste 16 autres billets ; $72+160=232$. La solution est entière et comprise entre 0 et 28.
\end{corrige}\end{blocLivre}
\begin{blocLivre}{cor-5-12}\begin{corrige}{12}{Un rectangle à retrouver}\demarche{Traduire le périmètre avec les deux côtés.}
Notons $x$ sa largeur, sa longueur vaut $x+7$. Alors $2x+2(x+7)=58$, donc $4x=44$ et $x=11$. Les dimensions sont $\res{11\text{ cm et }18\text{ cm}}$, l’aire $\res{198\text{ cm}^2}$.
\label{fig-cor-5-12}\figureGeo[corrige]{rectangle-equation}
\end{corrige}\end{blocLivre}
\begin{blocLivre}{cor-5-13}\begin{corrige}{13}{Un achat rentable}\demarche{Résoudre l’inéquation, puis revenir aux valeurs entières du contexte.}
L’inéquation est $38+12n\le30n$, soit $38\le18n$, donc $n\ge38/18\approx2{,}1$. Le nombre d’années étant entier, la gourde devient rentable à partir de la $\res{3^{\text{e}}\text{ année}}$. Contrôle : 74 € contre 90 € à trois ans, 62 € contre 60 € à deux ans.
\end{corrige}\end{blocLivre}
\begin{blocLivre}{cor-5-14}\begin{corrige}{14}{Deux aires égales}\demarche{Développer pour faire disparaître les termes carrés communs.}
On pose $(x+2)^2=x(x+5)$. Après développement : $x^2+4x+4=x^2+5x$, donc $\res{x=4}$. Le carré a pour côté 6 et le rectangle pour côtés 4 et 9. L’aire commune est $\res{36\text{ cm}^2}$. Toutes les dimensions sont positives.
\label{fig-cor-5-14}\figureGeo[corrige]{aires-egales}
\end{corrige}\end{blocLivre}
\begin{blocLivre}{cor-5-15}\begin{corrige}{15}{Analyse et synthèse}\demarche{Déduire les valeurs possibles, puis vérifier chacune dans le problème.}
On développe : $(x+1)^2-25=x^2+2x-24=x(x+2)-24$. Analyse : si $x$ convient, $(x+1)^2-5^2=0$, donc $(x-4)(x+6)=0$, soit $x=4$ ou $x=-6$. Synthèse : une largeur est positive, donc seul 4 peut convenir ; un rectangle de $\res{4\text{ cm sur }6\text{ cm}}$ a bien une aire de 24 cm$^2$.
\end{corrige}\end{blocLivre}
